In recent years, the invasive Asian hornet (\textit{Vespa velutina}) has appeared across Europe, posing a significant threat to local wildlife and apiaries. Multiple methods for containing this threat have been explored, including visual tracking, triangulation, harmonic radar, and radio telemetry.

This thesis investigates radio telemetry to track Asian hornets. The primary focus is to design a cost-effective, rechargeable, sub-$200\text{ mg}$ transmitter that is practical and user-friendly for non-technical operators. The user must be able to attach and activate the transmitter in a simple manner and track the hornet with a portable receiver. This thesis aims to improve upon existing wildlife trackers by being purpose-built specifically for the tracking and eradication of Asian hornets.

The transmitter must minimize energy consumption and component count. This is achieved using a simple, single-stage transmitter designed around a Pierce oscillator combined with a pulsing circuit to conserve power.

Weight minimization is a critical aspect of this research, offering key practical advantages: reduced tag weight allows a larger battery to be integrated to extend operational lifespan or increase RF output power, while enabling a higher proportion of hornets to carry the payload over greater distances.

The primary goals of this thesis are first to design, simulate, and fabricate an initial transmitter tag, which will be evaluated and refined into an optimized second design iteration. Once a functional tag is realized, a complementary receiver will be designed and tested, leading to field deployment once weight targets are met.



%This (short) summary, accessible to a wide audience, is written in English and contains a \textbf{maximum of 3500 characters}. This summary must also be uploaded to KU Loket.


Keywords: Vespa velutina, radio telemetry, Pierce oscillator %Enter (more or less) five keywords.
